\clearpage\pagestyle{otMain}\setlength{\otBlocEspacement}{2.5mm}\fontsize{11}{15}\selectfont\corrigesChapitre{10}{Inconnues et équations simples}
\begin{blocLivre}{cor-10-1}\begin{corrige}{1}{Une addition à trou}\demarche{Soustraire le terme connu.}
On retire 8 : $x=21-8=\res{13}$. Contrôle : $13+8=21$.
\end{corrige}\end{blocLivre}
\begin{blocLivre}{cor-10-2}\begin{corrige}{2}{Une multiplication à trou}\demarche{Diviser par le facteur connu.}
On divise par 6 : $x=42:6=\res7$. Contrôle : $6\times7=42$.
\end{corrige}\end{blocLivre}
\begin{blocLivre}{cor-10-3}\begin{corrige}{3}{Un quotient exact}\demarche{Conserver le quotient exact.}
$x=7:4=\res{7/4=1{,}75}$. Contrôle : $4\times(7/4)=7$.
\end{corrige}\end{blocLivre}
\begin{blocLivre}{cor-10-4}\begin{corrige}{4}{Tester un candidat}\demarche{Substituer dans l’équation.}
$3\times5=15\ne18$. Donc 5 n’est pas une solution. La solution est $18:3=\res6$.
\end{corrige}\end{blocLivre}
\begin{blocLivre}{cor-10-5}\begin{corrige}{5}{Une valeur négative}\demarche{Autoriser les nombres relatifs.}
$x=4-9=\res{-5}$. Contrôle : $-5+9=4$.
\end{corrige}\end{blocLivre}
\begin{blocLivre}{cor-10-6}\begin{corrige}{6}{Avec un décimal}\demarche{Aligner les rangs décimaux.}
$x=8{,}15-2{,}70=\res{5{,}45}$. On contrôle par l’addition.
\end{corrige}\end{blocLivre}
\begin{blocLivre}{cor-10-7}\begin{corrige}{7}{Un coefficient décimal}\demarche{Diviser par le coefficient non nul.}
$x=6:0{,}8=60:8=\res{7{,}5}$. Contrôle : $0{,}8\times7{,}5=6$.
\end{corrige}\end{blocLivre}
\begin{blocLivre}{cor-10-8}\begin{corrige}{8}{Additionner un négatif}\demarche{Ajouter l’opposé du terme connu.}
$x=2-(-6)=2+6=\res8$. Contrôle : $8+(-6)=2$.
\end{corrige}\end{blocLivre}
\begin{blocLivre}{cor-10-9}\begin{corrige}{9}{Trois équations proches}\demarche{Distinguer terme, facteur et diviseur.}
La première donne $x=12-7=\res5$. La deuxième donne $x=12:7=\res{12/7}$. Pour la troisième, on inverse la division par 7 : $x=12\times7=\res{84}$. L’opération inverse dépend de l’opération présente.
\end{corrige}\end{blocLivre}
\begin{blocLivre}{cor-10-10}\begin{corrige}{10}{Une mauvaise opération}\demarche{Identifier l’opération à inverser.}
Il faut inverser une multiplication, donc diviser : $x=23:5=\res{4{,}6}$. Contrôle : $5\times4{,}6=23$. $5\times18$ ne vaut pas 23.
\end{corrige}\end{blocLivre}
\begin{blocLivre}{cor-10-11}\begin{corrige}{11}{Un prix unitaire}\demarche{Nommer l’inconnue et son unité.}
Si $x$ est le prix d’un cahier en euros, $9x=25{,}20$. Donc $x=25{,}20:9=\res{2{,}80}$ euros. Le contrôle donne bien le prix de neuf cahiers.
\end{corrige}\end{blocLivre}
\begin{blocLivre}{cor-10-12}\begin{corrige}{12}{Une quantité initiale}\demarche{Traduire l’évolution donnée.}
Si $x$ est le volume initial en litres, $x+17=63$. Donc $x=63-17=\res{46\ \mathrm L}$.
\end{corrige}\end{blocLivre}
\begin{blocLivre}{cor-10-13}\begin{corrige}{13}{Aire et largeur}\demarche{Utiliser la formule de l’aire.}
Pour la largeur $x$ en cm, $8x=46$. Ainsi $x=46:8=\res{5{,}75\ \mathrm{cm}}$. Le résultat est une longueur, même si la donnée est une aire.
\label{fig-cor-10-13}\figureGeo[corrige]{eq-13}
\end{corrige}\end{blocLivre}
\begin{blocLivre}{cor-10-14}\begin{corrige}{14}{Parts dans un schéma}\demarche{Isoler les parts avant de partager.}
On retire d’abord le supplément : $45-6=39$ euros pour les trois parts. Une part vaut $39:3=\res{13}$ euros. Contrôle : $13+13+13+6=45$.
\label{fig-cor-10-14}\figureGeo[corrige]{eq-14}
\end{corrige}\end{blocLivre}
\begin{blocLivre}{cor-10-15}\begin{corrige}{15}{Une contrainte d’objets entiers}\demarche{Confronter la solution aux contraintes du contexte.}
Numériquement, $x=50/8=\res{6{,}25}$ résout l’équation. Mais une part de 6,25 jetons n’est pas admise. En objets entiers, le partage égal sans reste est impossible : $50=8\times6+2$.
\end{corrige}\end{blocLivre}
